\documentclass[../RFC-HDFG-2026-001.tex]{subfiles}

\begin{document}

%% ======================================================================
%%  Section 3 – Parameter String Format
%% ======================================================================
\section{Parameter String Format}
\label{sec:param-format}

The parameter string format is based on TOML v1.0.0~\cite{toml10}, parsed by
a vendored subset of the tomlc17 library~\cite{tomlc17}.  The full TOML
v1.0.0 specification is the normative reference for syntax and value types;
this section documents the constraints imposed by the library and the
rationale for deviations.  Two usage profiles are defined: a tighter
\textbf{filter profile} for I/O filter configuration strings, and a more
permissive \textbf{VFD/VOL profile} for Virtual File Driver and VOL connector
configuration strings.

Using an established, externally maintained grammar provides typed values
(integer, float, boolean, string, array, inline table), a stable specification
the library does not need to own, and a path to nested configuration for VOL
connector and VFD stacks without changing the API or the wire format.

\subsection{Format Overview}
\label{sec:param-overview}

A parameter string is a comma-separated list of \texttt{key = value} pairs
conforming to the TOML inline table content grammar.  Outer braces are
\emph{optional}: the library accepts both the bare form
(\texttt{level = 6}) and the braced form (\texttt{\{level = 6\}}), wrapping
the string in \texttt{\{\ldots\}} internally before parsing if braces are
absent.  \texttt{NULL} or an empty string means ``no parameters'' and is
valid for filters that take no parameters (shuffle, fletcher32, nbit).

\verysubsection{Value types}

\begin{itemize}
\item \textbf{Integer:} Unquoted decimal integer, optional leading sign.
  Example: \texttt{level = 6}, \texttt{offset = -4}.
  TOML integer syntax (underscores as digit separators, \texttt{0x}/\texttt{0o}/\texttt{0b}
  prefixes) is accepted by the parser but filter callbacks receive the
  parsed \texttt{int64\_t} value directly via \texttt{H5Zconfig\_get\_int}.
\item \textbf{Float:} Unquoted decimal float with mandatory decimal point or
  exponent. Example: \texttt{rate = 3.5}, \texttt{tol = 1.0e-6}.
  The decimal separator is always a period regardless of locale.
  TOML special float values (\texttt{inf}, \texttt{nan}) are accepted by the
  parser but rejected with \texttt{H5E\_BADVALUE} by the library before
  \texttt{set\_config} is invoked.
\item \textbf{Boolean:} \texttt{true} or \texttt{false} (lowercase only, per
  TOML).  The values \texttt{True}, \texttt{False}, \texttt{1}, and \texttt{0}
  are \emph{not} recognized as booleans; the parser rejects them with
  \texttt{H5E\_BADVALUE}.  Example: \texttt{fast\_mode = true}.
\item \textbf{String:} Double-quoted only (\texttt{"\ldots"}).  Standard TOML
  escape sequences are supported (\texttt{\textbackslash n}, \texttt{\textbackslash t},
  \texttt{\textbackslash\textbackslash}, \texttt{\textbackslash"}, etc.).
  Example: \texttt{coding = "entropy"}, \texttt{path = "/data/run\_1"}.
  Single-quoted TOML literal strings (\texttt{'\ldots'}) are also accepted;
  no escape processing is performed inside single quotes.
  An empty quoted value (\texttt{key = ""} or \texttt{key = ''}) is valid
  and passes the empty string to the filter callback; it is the filter's
  responsibility to accept or reject an empty value.
  A bare \texttt{key=} with no value and no quotes is a parse error and is
  rejected with \texttt{H5E\_BADVALUE}.
\item \textbf{Inline table (nested):} A value may itself be a
  comma-separated list of key=value pairs enclosed in braces.
  Example: \texttt{stripe = \{size = 1048576, count = 4\}}.
  Nesting is appropriate for VOL connector and VFD configuration, and
  also for filters that wrap a sub-compressor with its own parameters
  (e.g., Blosc selecting an internal compressor:
  \texttt{compressor = \{name = "zlib", level = 6\}, shuffle = 1}).
  Dotted-key form (\texttt{compressor.name = "zlib", compressor.level = 6})
  is equivalent and is the recommended style for readability in a
  single-line string.
\end{itemize}

\verysubsection{Keys}

Keys follow the TOML bare-key rules: ASCII letters, digits, hyphens, and
underscores (\texttt{[A-Za-z0-9\_-]+}).  Quoted keys (\texttt{"key"}) are
accepted by the parser but discouraged.  Keys are \emph{case-insensitive};
the library normalizes to lowercase before lookup.  Dotted keys
(\texttt{stripe.size}) are accepted and express one level of nesting
(\texttt{stripe.size = 1048576} is equivalent to
\texttt{stripe = \{size = 1048576\}}); they are the recommended form for
VOL/VFD nested parameters as they are more readable in a single-line string.
Duplicate keys are rejected with \texttt{H5E\_BADVALUE}.

\verysubsection{Additional constraints — filter profile}

The following constraints apply to parameter strings used with I/O filters
(i.e., strings passed to \texttt{H5Pappend\_filter} and the associated
\texttt{set\_config}/\texttt{get\_config} callbacks).

\begin{itemize}
\item \textbf{Reserved key:} The key \texttt{filter\_title} is reserved by the
  library (\SectionRef{sec:class3-struct}) and rejected with
  \texttt{H5E\_BADVALUE} if it appears in a parameter string.
\item \textbf{Maximum string length:} \texttt{H5Z\_CONFIG\_STRING\_MAX} (= 4096
  bytes).  Strings exceeding this limit are rejected with
  \texttt{H5E\_BADVALUE}.
\item \textbf{Maximum parameter count:} \texttt{H5Z\_CONFIG\_MAX\_PARAMS} (= 64
  top-level keys).  Strings containing more top-level keys are rejected with
  \texttt{H5E\_BADVALUE}.  This bound limits parsing work and exposure to
  maliciously crafted strings.
\item \textbf{Special float values:} \texttt{inf}, \texttt{-inf}, and
  \texttt{nan} are valid TOML but rejected by the library with
  \texttt{H5E\_BADVALUE}.
\item \textbf{TOML features not supported:} Arrays (\texttt{[1,2,3]}),
  multi-line strings, datetime values, and document-level tables
  (\texttt{[section]}) are not part of the filter profile.  The parser
  rejects them with \texttt{H5E\_BADVALUE}.
\item \textbf{Locale independence:} The parser and all typed accessors use
  C-locale numeric conversion.  Process-wide locale settings do not affect
  parsing.  The decimal separator is always a period.
\end{itemize}

\verysubsection{Additional constraints — VFD/VOL profile}

VFD and VOL connector configuration strings require richer structure than
filter parameter strings: arbitrary stack depth, arrays (e.g., lists of
endpoint addresses), and file-based delivery that benefits from multi-section
TOML documents.  The following relaxed constraints apply to parameter strings
used in VFD and VOL connector configuration contexts.

\begin{itemize}
\item \textbf{Full TOML v1.0.0 document syntax supported:} Arrays
  (\texttt{[1,2,3]}), multi-line strings, and document-level tables
  (\texttt{[section]}) are all accepted.  The \texttt{filter\_title} reserved
  key does not apply in this profile.
\item \textbf{Arbitrary nesting depth:} Inline tables and dotted keys may be
  nested to any depth required to express the VFD stack hierarchy.  There is
  no library-imposed depth limit beyond available stack memory.
\item \textbf{Maximum string length:} \texttt{H5FD\_CONFIG\_STRING\_MAX} (=
  16384 bytes).  A full multi-level VFD stack with a SWMR configuration group
  is unlikely to exceed a few kilobytes; 16\,KB provides ample headroom.
  Strings exceeding this limit are rejected with \texttt{H5E\_BADVALUE}.
\item \textbf{Maximum top-level key count:} \texttt{H5FD\_CONFIG\_MAX\_PARAMS}
  (= 256).  Strings containing more top-level keys are rejected with
  \texttt{H5E\_BADVALUE}.
\item \textbf{Special float values:} \texttt{inf}, \texttt{-inf}, and
  \texttt{nan} are rejected with \texttt{H5E\_BADVALUE}, consistent with the
  filter profile.
\item \textbf{Locale independence:} C-locale numeric conversion applies, as
  in the filter profile.
\end{itemize}

The VFD name is expressed as the sole top-level key, with its configuration
parameters as the value (an inline table or dotted keys).  An underlying VFD
in a stack is expressed as a nested key under the outer VFD's table.  An
empty parameter list is expressed as an empty inline table (\texttt{\{\}}).
For example, a three-level page-buffer~$\to$~encryption~$\to$~sec2 stack:

\begin{minted}{toml}
page_buffer.page_size       = 4096
page_buffer.max_num_pages   = 64
page_buffer.underlying_vfd.encryption_vfd.plaintext_page_size = 4096
page_buffer.underlying_vfd.encryption_vfd.key = "0123456789ABCDEF..."
page_buffer.underlying_vfd.encryption_vfd.underlying_vfd.sec2 = {}
\end{minted}

or equivalently as inline tables:

\begin{minted}{toml}
page_buffer = {
  page_size = 4096, max_num_pages = 64,
  underlying_vfd = { encryption_vfd = {
    plaintext_page_size = 4096, key = "0123456789ABCDEF...",
    underlying_vfd = { sec2 = {} }
  }}
}
\end{minted}

\subsection{VFD/VOL Profile: Use Case Coverage}
\label{sec:vfd-use-cases}

This subsection demonstrates, use case by use case, that the TOML VFD/VOL
profile is sufficient to express all VFD stack configuration requirements.
The Lifeboat VFD configuration RFC~\cite{lifeboat-vfd-cl} provides a concrete
set of use cases and motivates several of the design choices described here;
it is used as a reference point, not a competing proposal.

\verysubsection{UC-1: Flat scalar VFD parameters}

\textbf{Use case.}  A VFD such as the page buffer VFD requires a small set of
scalar parameters: page size, maximum number of pages, replacement policy.

\textbf{Coverage.}  Standard TOML key\,=\,value pairs cover this
directly.  The grammar provides typed integers, floats, booleans, and quoted
strings without any extension.  Example:

\begin{minted}{toml}
page_buffer = { page_size = 4096, max_num_pages = 64, replacement_policy = 0 }
\end{minted}

No aspect of this use case requires any relaxation beyond ordinary TOML.

\verysubsection{UC-2: Boolean flags}

\textbf{Use case.}  VFD SWMR configuration contains several boolean fields
(e.g., \texttt{presume\_posix\_semantics}, \texttt{flush\_raw\_data}).  The
Lifeboat grammar has no boolean type; it uses integer \texttt{0}/\texttt{1}
as a workaround, requiring comment annotations (\texttt{/* <- Use integer for
bool value */}) to explain the intent.

\textbf{Coverage.}  TOML has a native boolean type:
\texttt{true}/\texttt{false} (lowercase, unquoted).  The intent is
self-documenting and the parser enforces the type without integer-to-boolean
coercion.  Example:

\begin{minted}{toml}
H5F_vfd_swmr_config = {
  presume_posix_semantics = true,
  flush_raw_data          = true,
  generate_updater_files  = false
}
\end{minted}

Native boolean support eliminates the need for integer \texttt{0}/\texttt{1} workarounds.

\verysubsection{UC-3: Binary blob values (e.g., encryption keys)}

\textbf{Use case.}  The encryption VFD requires a raw binary key (32 bytes in
the Lifeboat example).  The Lifeboat grammar introduces a special
\texttt{<binary\_blob>} terminal with a \texttt{--} prefix to distinguish hex
byte sequences from identifiers and integers.

\textbf{Coverage.}  TOML has no native binary type, but a quoted
string carrying a hex-encoded value is fully sufficient.  The VFD callback
receives the string and decodes it.  The convention is explicit and requires no
grammar extension:

\begin{minted}{toml}
encryption_vfd = {
  key = "0123456789ABCDEF0123456789ABCDEF0123456789ABCDEF0123456789ABCDEF"
}
\end{minted}

TOML's quoted string type carries arbitrary content, so no disambiguation
prefix is needed.  The decoding convention (e.g., hex, base64) is specified
by the VFD and validated in its configuration entry point, as with all other
parameter types.

\textbf{Validation responsibility.}  In the Lifeboat grammar, the parser
natively enforces that a \texttt{<binary\_blob>} terminal contains only valid
hex characters.  Under TOML the parser treats the value as an opaque string;
it is therefore the VFD's responsibility to validate the string's content
(character set, byte length, encoding) and return \texttt{H5E\_BADVALUE} if
the value is malformed.  This is the same responsibility pattern already
required of all other parameter types: the TOML grammar validates syntax and
type, while semantic validation (range checks, format checks) belongs to the
VFD.  The shift is in degree, not in kind.

\verysubsection{UC-4: Arbitrary-depth VFD stack nesting}

\textbf{Use case.}  A VFD stack may contain three or more levels (e.g.,
page buffer~$\to$~encryption~$\to$~sec2).  Each VFD must be able to read its
own parameters and forward the remaining configuration to the underlying VFD,
without any VFD needing to know the schema of its neighbors.

\textbf{Coverage.}  TOML inline tables and dotted keys are
recursively defined and support arbitrary nesting depth.  The VFD/VOL profile
removes any library-imposed depth limit.  Each VFD in the stack reads the
keys it recognizes by name and extracts the \texttt{underlying\_vfd} sub-table
to pass to the next level.  Because TOML keys are named rather than
positional, each VFD looks up only what it knows about and ignores everything
else; no VFD needs visibility into any other VFD's schema.

The full three-level example from UC-1:

\begin{minted}{toml}
page_buffer = {
  page_size = 4096, max_num_pages = 64,
  underlying_vfd = { encryption_vfd = {
    plaintext_page_size = 4096, key = "...",
    underlying_vfd = { sec2 = {} }
  }}
}
\end{minted}

The page buffer VFD reads \texttt{page\_size} and \texttt{max\_num\_pages},
then calls \texttt{H5FDconfig\_get\_subtable\_str} to extract the
\texttt{underlying\_vfd} value as a raw TOML string
(\SectionRef{sec:subtable-extraction}), creates a new FAPL, and loads that
string into the FAPL's \texttt{driver\_config\_str} field via
\texttt{H5FDconfig\_load\_str\_into\_fapl}.  The encryption VFD
receives that FAPL, reads its own keys from the extracted string, and
similarly forwards the \texttt{underlying\_vfd.sec2} sub-table to the sec2
configuration.  No VFD ever parses keys belonging to another VFD.

\textbf{API note.}  VFDs do not use the \texttt{H5Z\_class3\_t.set\_config}
callback defined in \SectionRef{sec:class3}; that callback is specific to I/O
filters.  VFD configuration strings are delivered via the
\texttt{driver\_config\_str} field of \texttt{H5FD\_driver\_prop\_t} (defined
in \texttt{H5FDprivate.h}), which is already present in the HDF5 library.
The TOML VFD/VOL profile described here specifies the \emph{format} of that
string; no change to the \texttt{H5FD\_class\_t} struct or the VFD open
callback signature is required by this RFC.

\verysubsection{UC-5: VFD name as the top-level identifier}

\textbf{Use case.}  In the Lifeboat grammar, the VFD name is the outer
identifier of the top-level name-value pair: \texttt{( page\_buffer ( \ldots )
)}.  The configuration string is self-describing: the consumer can determine
which VFD is being configured without out-of-band information.

\textbf{Coverage.}  The same self-describing property is achieved
by convention: the VFD name is the sole top-level key, and its value is the
parameter table.  A configuration string that begins with
\texttt{page\_buffer~=~\{~\ldots~\}} unambiguously identifies the page buffer
VFD as the top-level driver.  The library reads the top-level key name to
determine which VFD to load and configure.

\textbf{Sole top-level key enforcement.}  The library validates that a
VFD/VOL profile configuration string contains \emph{exactly one} top-level
key.  If zero top-level keys are present the string is treated as ``no
configuration'' (equivalent to a \texttt{NULL} string).  If two or more
top-level keys are present (e.g., \texttt{page\_buffer = \{\ldots\}, sec2 =
\{\}}) the library returns \texttt{H5E\_BADVALUE} immediately, before
consulting any VFD.  This rule ensures unambiguous driver identification and
mirrors the Lifeboat grammar's constraint that the top-level is always a
single \texttt{<name\_value\_pair>}.

\verysubsection{UC-6: Empty parameter list}

\textbf{Use case.}  The sec2 VFD requires no configuration parameters.  In
the Lifeboat grammar this is expressed as \texttt{( sec2 () )} --- a VFD name
with an explicit but empty parameter list.

\textbf{Coverage.}  An empty inline table \texttt{sec2 = \{\}} is
the direct and idiomatic equivalent.  The VFD name is present (self-describing),
the parameter list is explicitly empty, and the grammar requires no special
case.  \texttt{NULL} or an empty string is also accepted by the VFD/VOL
profile and treated as ``no parameters'', consistent with the filter profile.

\verysubsection{UC-7: Configuration groups (VFD SWMR)}

\textbf{Use case.}  VFD SWMR requires coordinated configuration across
multiple subsystems: the SWMR engine itself, the page buffer, file space
strategy, and file space page size.  Historically these were configured
independently, risking writer/reader mismatches.  The Lifeboat RFC introduces
a ``configuration group'' concept: a named outer wrapper containing multiple
named sub-configurations that are parsed and applied together.

\textbf{Coverage.}  The configuration group maps directly onto a
TOML table whose value is a list of named sub-tables.  When delivered as an
inline string (e.g., environment variable or API call), nested inline tables
express the group:

\begin{minted}{toml}
vfd_swmr_config_data = {
  H5F_vfd_swmr_config = {
    version = 1, tick_len = 4, max_lag = 7,
    presume_posix_semantics = true, flush_raw_data = true,
    md_file_path = "./md_dir", md_file_name = "md_file"
  },
  page_buffer_config = {
    page_buf_size = 4096, metadata_pages_only = true
  },
  file_space_strategy_config = { persist = false },
  file_space_page_size        = { page_size = 4096 }
}
\end{minted}

When delivered from a file, standard TOML document-level section tables are
the natural form (permitted under the VFD/VOL profile):

\begin{minted}{toml}
[vfd_swmr_config_data.H5F_vfd_swmr_config]
version   = 1
tick_len  = 4
max_lag   = 7
presume_posix_semantics = true
flush_raw_data          = true
md_file_path = "./md_dir"
md_file_name = "md_file"

[vfd_swmr_config_data.page_buffer_config]
page_buf_size       = 4096
metadata_pages_only = true

[vfd_swmr_config_data.file_space_strategy_config]
persist = false

[vfd_swmr_config_data.file_space_page_size]
page_size = 4096
\end{minted}

Both forms parse to the same TOML AST and express identical configuration
data.  However, they are not processed identically at the library level: an
inline string is stored verbatim in the FAPL's \texttt{driver\_config\_str}
and consumed directly by VFDs, while a file containing document-level section
tables is first parsed and normalized to an inline string by the
file-loading helper (\SectionRef{sec:subtable-extraction}) before being
placed in the FAPL.  VFDs always see the inline form; the distinction between
the two authoring styles is transparent to them.

The top-level name (\texttt{vfd\_swmr\_config\_data}) identifies the
configuration group, and each sub-table name identifies the target subsystem.
The writer/reader coordination
problem is solved by the same mechanism: both sides load from the same
configuration string or file, with the library applying writer- and
reader-specific flag overrides (\texttt{writer}, \texttt{create\_file})
programmatically rather than encoding them in the configuration string.

\verysubsection{UC-8: Stack-level delegation without breadth-first parsing}

\textbf{Use case.}  The Lifeboat grammar was designed around a
\emph{breadth-first} parsing model: when a VFD encounters the configuration
string, it parses only its own parameters and receives the sub-expression for
the underlying VFD as a single opaque token, which it passes down the stack
unchanged.  The motivation, as stated in the Lifeboat RFC, is schema
isolation: each VFD in the stack should not need to know anything about the
configuration data for the underlying VFDs.  This is achieved by a
context-sensitive lexer that, when a name-value pair value is expected,
treats a parenthesised sub-expression as a single token rather than
descending into it.  Each VFD parses only the keys at its own level and
forwards the unparsed sub-expression string down the stack.

\textbf{Schema isolation via AST lookup.}  TOML parsers (including the
vendored tomlc17 subset) build a full recursive AST in a single pass, after
which each VFD can locate the keys it recognises by name and extract the
\texttt{underlying\_vfd} sub-table without scanning past unknown structure.
Schema isolation is preserved --- each VFD reads only its own keys --- through
named-key lookup over a complete AST.

\textbf{Implementation.}  Each VFD in the stack receives a TOML configuration
string via \texttt{driver\_config\_str} in its FAPL.  It parses that string,
extracts its own keys by name using \texttt{H5FDconfig\_get\_*} helpers, then
calls \texttt{H5FDconfig\_get\_subtable\_str} (\SectionRef{sec:subtable-extraction})
to extract the raw source bytes of the \texttt{underlying\_vfd} value — no
AST-to-string serialization is required, because the helper works directly on
the original input bytes.  The extracted substring is loaded into a new FAPL's
\texttt{driver\_config\_str} and passed to the next VFD's open callback.
Configuration strings are small (well under the 16\,KB VFD/VOL profile limit),
so the substring extraction overhead is negligible.  Each VFD independently
validates and consumes only its own configuration data, achieving full schema
isolation without any special parsing protocol.

\subsection{Examples}
\label{sec:param-examples}

\begin{center}
\begin{tabularx}{\textwidth}{>{\ttfamily\raggedright\arraybackslash\hsize=1.55\hsize}X
                             >{\raggedright\arraybackslash\hsize=0.45\hsize}X}
\toprule
\textbf{Parameter String} & \textbf{Use Case} \\
\midrule
level = 6 & deflate (integer) \\
mode = "rate", rate = 3.5 & ZFP (string + float) \\
pixels\_per\_block = 16, coding = "entropy" & SZIP \\
scale\_type = "float\_dscale", scale\_factor = 4 & scaleoffset \\
fast\_mode = true, level = 6 & boolean flag \\
path = "/data/run\_1,\allowbreak{}v2/dict.bin" & string value with a comma \\
compressor.name = "zlib",\allowbreak{} compressor.level = 6,\allowbreak{} shuffle = 1
  & Blosc2: dotted keys \\
compressor = \{name = "zlib",\allowbreak{} level = 6\},\allowbreak{} shuffle = 1
  & Blosc2: inline table (equivalent) \\
stripe = \{size = 1048576, count = 4\} & nested table (VOL/VFD) \\
NULL & no parameters \\
\bottomrule
\end{tabularx}
\end{center}

\subsection{Parser Implementation}
\label{sec:param-parser}

The parameter string parser is a vendored, reduced subset of
tomlc17~\cite{tomlc17}, a MIT-licensed C17 TOML library.  The vendored
copy lives under \texttt{src/H5Ztoml.\{c,h\}} and is stripped to the inline
table and primitive value subset described in \SectionRef{sec:param-overview}.
HDF5's memory allocation functions (\texttt{H5MM\_malloc} etc.)\ are
substituted for the library's default allocator.  The vendored copy is
updated as part of the normal HDF5 dependency refresh process; the TOML
v1.0.0 grammar is stable and no API-breaking changes are expected.

The semicolon character is reserved at the grammar level for a possible
future pipeline-string extension (\texttt{"shuffle; deflate, level = 9;
fletcher32"}); it is not part of the TOML inline table grammar and is
rejected with \texttt{H5E\_BADVALUE} outside of quoted strings.

\verysubsection{Sub-table string extraction for VFD stack delegation}
\label{sec:subtable-extraction}

VFD stack delegation (\SectionRef{sec:vfd-use-cases}, UC-4 and UC-8) requires
a VFD to pass the configuration of its underlying VFD as a string to that
VFD's configuration entry point.  Because the \texttt{set\_config} entry point
accepts a \texttt{const char~*} string, a VFD cannot pass a parsed
\texttt{toml\_table\_t} pointer directly to the next level.

Standard C TOML parsers, including tomlc17, build an in-memory parse tree but
do not include an emitter to serialize that tree back to a string.  Rather
than adding a full TOML emitter to the vendored library---which would
significantly increase complexity and maintenance burden---the library instead
provides a lightweight helper function that extracts the raw source substring
corresponding to a named sub-table directly from the original input string:

\begin{minted}{c}
herr_t H5FDconfig_get_subtable_str(
    const char  *params,      /* original VFD config string          */
    const char  *key,         /* dotted key path, e.g. "underlying_vfd" */
    char        *buf,         /* output buffer, or NULL for size query   */
    size_t      *buf_size     /* in: capacity; out: required length      */
);
\end{minted}

This function locates the value span of \texttt{key} in \texttt{params} by
scanning the source bytes directly (using the same lexer used by the parser
to identify balanced brace extents), copies the raw substring into
\texttt{buf}, and returns its length via \texttt{buf\_size}.  The extracted
substring is a valid TOML inline table string that the receiving VFD can
parse independently.  This approach:

\begin{itemize}
\item Requires no emitter and no AST traversal beyond balanced-brace matching.
\item Preserves the exact source text, including whitespace and comments,
  avoiding any round-trip fidelity issues.
\item Follows the standard HDF5 size-query pattern: passing \texttt{NULL} for
  \texttt{buf} returns the required buffer size without copying.
\item Returns \texttt{H5E\_NOTFOUND} if the key is absent, and
  \texttt{H5E\_BADVALUE} if the value is not a table (i.e., is a scalar).
\end{itemize}

\textbf{Inline table requirement.}
\texttt{H5FDconfig\_get\_subtable\_str} operates on inline table syntax only.
Balanced-brace scanning works because the value of a named key in an inline
table is always a contiguous, brace-delimited span in the source bytes.
Document-level section tables (\texttt{[section]}) have no such span: the
keys of \texttt{[page\_buffer.underlying\_vfd]} are scattered across the file
with no enclosing braces, so there is no contiguous substring to extract.

This is not a limitation in practice because document-level tables
\emph{never reach \texttt{driver\_config\_str} directly}.  The two delivery
paths are:

\begin{enumerate}
\item \textbf{Inline string} (environment variable, API call): the caller
  supplies an inline table string, which is stored verbatim in
  \texttt{driver\_config\_str}.  \texttt{H5FDconfig\_get\_subtable\_str}
  operates on this string directly.
\item \textbf{File-based delivery}: the application calls a library helper
  such as \texttt{H5FDconfig\_load\_from\_file()} (analogous to the Lifeboat
  \texttt{H5CL\_load\_config\_string\_from\_file()}), which reads the file,
  parses it into a TOML AST (accepting document-level tables in that parse),
  and then emits a single normalized inline table string that it stores in
  the FAPL's \texttt{driver\_config\_str}.  From that point on, all VFDs in
  the stack see only the inline form.
\end{enumerate}

The document-level \texttt{[section]} syntax is therefore a user-facing
authoring convenience.  The library's file-loading helper is the single point
responsible for normalizing it to inline form; \texttt{H5FDconfig\_get\_subtable\_str}
and all downstream VFD configuration logic operate exclusively on inline
strings.  The normalization step requires the file-loading helper to include a
minimal inline-table emitter, scoped only to that one helper and not exposed
to VFD authors.

This function is defined in \texttt{H5FDconfig.c} and declared in
\texttt{H5FDdevelop.h}, placing it in the VFD developer API rather than the
filter developer API (\texttt{H5Zdevelop.h}).  The \texttt{H5FD} prefix
distinguishes it from the filter-scoped \texttt{H5Zconfig\_*} helpers.


%% ======================================================================
%%  Section 4 – Filter Plugin Class Version 3
%% ======================================================================
\section{Filter Plugin Class Version 3}
\label{sec:class3}

\subsection{Rationale}

The string configuration feature requires filters to participate voluntarily.
A filter must be able to accept a parameter string and produce \texttt{cd\_values},
and optionally to reverse the process for introspection. This is achieved by
introducing two new optional callbacks in a version~3 filter class structure.

\subsection{Version Numbering}
\label{sec:version-numbering}

The existing version-to-struct mapping in HDF5 is:

\begin{center}
\begin{tabular}{ll}
\toprule
\textbf{Struct} & \textbf{Version Field} \\
\midrule
\texttt{H5Z\_class1\_t} & none (legacy) \\
\texttt{H5Z\_class2\_t} & \texttt{1} \\
\texttt{H5Z\_class3\_t} & \texttt{2} \\
\bottomrule
\end{tabular}
\end{center}

The version field is a single, monotonically increasing integer shared
across all class struct generations.  Plugin authors set this field to the
literal value corresponding to their chosen struct (\texttt{1} for
\texttt{H5Z\_class2\_t}, \texttt{2} for \texttt{H5Z\_class3\_t}).  Adding a
future \texttt{H5Z\_class4\_t} would assign it value \texttt{3}, with no new
per-struct version macro required.

The version constants are:

\begin{minted}{c}
#define H5Z_CLASS_T_VERS      1   /* Unchanged -- existing constant for H5Z_class2_t  */
#define H5Z_CLASS_T_VERS_MAX  2   /* Highest accepted version field value (currently  *
                                   * H5Z_class3_t).  Used by the library for sanity   *
                                   * checking; plugin authors should not reference it *
                                   * directly, since its value increases as new       *
                                   * struct generations are added.                    */
\end{minted}

Note that \texttt{H5Z\_CLASS\_T\_VERS} remains at \texttt{1} for backward
compatibility with code that already references it.  No new
\texttt{H5Z\_CLASS3\_T\_VERS} (or future per-struct) constant is introduced;
naming a constant after a specific struct generation would not generalise
when \texttt{H5Z\_class4\_t} is added.  Plugin authors instead set the
\texttt{version} field to the literal value documented in the table above.

\textbf{Future class versions:} Other in-progress HDF5 work (e.g.,
Accelerator-native I/O Pipeline) may also need to extend
\texttt{H5Z\_class\_t}.  The \texttt{H5Z\_CLASS\_T\_VERS\_MAX} pattern
accommodates this without back-to-back struct versioning: a future
\texttt{H5Z\_class4\_t} would set \texttt{version = 3} and
\texttt{H5Z\_CLASS\_T\_VERS\_MAX} would be bumped to \texttt{3} in the same
release, with a new dispatch branch added for \texttt{version == 3}.
Extensions to \texttt{H5Z\_class\_t} from different features should be
merged in the same release to avoid consecutive struct-generation bumps.

The dispatch logic in \texttt{H5Zregister()} currently works by checking
\texttt{version != H5Z\_CLASS\_T\_VERS} first; any other value is treated as a
legacy \texttt{H5Z\_class1\_t} struct.  The dispatch is extended to recognise
\texttt{version == 2} as \texttt{H5Z\_class3\_t}, with
\texttt{H5Z\_CLASS\_T\_VERS\_MAX} guarding against unknown future values:

\begin{minted}{c}
    if (version > H5Z_CLASS_T_VERS_MAX)
        /* -> error: version field exceeds library support */
    else if (version == H5Z_CLASS_T_VERS)  /* 1 */
        /* -> H5Z_class2_t        */
    else if (version == 2)
        /* -> H5Z_class3_t        */
    else
        /* -> legacy H5Z_class1_t */
\end{minted}

Note: When \texttt{H5\_NO\_DEPRECATED\_SYMBOLS} is defined, the legacy
\texttt{H5Z\_class1\_t} path is compiled out entirely and any value other
than \texttt{H5Z\_CLASS\_T\_VERS} or \texttt{2} produces a hard error.  Under
\texttt{H5\_NO\_DEPRECATED\_SYMBOLS}, only versions \texttt{1} and \texttt{2}
are accepted.

\texttt{H5Z\_class3\_t} is an independent flat struct; it does NOT embed
\texttt{H5Z\_class2\_t} as a member.  The library reads the \texttt{version}
field in \texttt{H5Zregister()} via \texttt{memcpy}
through a \texttt{char~*} to avoid strict-aliasing undefined behaviour.
\texttt{H5PL\_\_open} receives an already-typed pointer via the plugin's
exported \texttt{H5PLget\_plugin\_info()} symbol and does not need this.

Document the version-to-struct mapping table in \texttt{H5Zdevelop.h}.

\verysubsection{Compatibility macro in \texttt{H5version.h}}
The existing compatibility macro
(\texttt{H5version.h}) maps the generic type alias \texttt{H5Z\_class\_t} to a versioned
struct:

\begin{minted}{c}
#if !defined(H5Z_class_t_vers) || H5Z_class_t_vers == 2
  #define H5Z_class_t H5Z_class2_t
#elif H5Z_class_t_vers == 1
  #define H5Z_class_t H5Z_class1_t
#endif
\end{minted}

Currently \texttt{H5Z\_class\_t\_vers} defaults to \texttt{2}, mapping
\texttt{H5Z\_class\_t} $\rightarrow$ \texttt{H5Z\_class2\_t}. This macro is not changed to default to
\texttt{3} (mapping to \texttt{H5Z\_class3\_t}), because existing code using \texttt{H5Z\_class\_t}
would silently get the larger struct with uninitialized callback fields. The
default remains \texttt{2}. A new value of \texttt{3} can be added as an opt-in for
code that explicitly requests \texttt{H5Z\_class3\_t}:

\begin{minted}{c}
#if !defined(H5Z_class_t_vers) || H5Z_class_t_vers == 2
  #define H5Z_class_t H5Z_class2_t
#elif H5Z_class_t_vers == 1
  #define H5Z_class_t H5Z_class1_t
#elif H5Z_class_t_vers == 3
  #define H5Z_class_t H5Z_class3_t
#endif
\end{minted}

Plugin authors adopting v3 should use \texttt{H5Z\_class3\_t} directly rather than
relying on the \texttt{H5Z\_class\_t} alias.


\subsection{Structure Definition}
\label{sec:class3-struct}

\begin{minted}{c}
typedef herr_t (*H5Z_set_config_func_t)(
    const char   *params,
    unsigned     *flags,
    size_t       *cd_nelmts,
    unsigned      cd_values[],
    size_t        cd_values_size
);

typedef herr_t (*H5Z_get_config_func_t)(
    unsigned          flags,
    size_t            cd_nelmts,
    const unsigned    cd_values[],
    char             *buf,
    size_t           *buf_size
);

typedef struct H5Z_class3_t {
    int                        version;
    H5Z_filter_t               id;
    unsigned                   encoder_present;
    unsigned                   decoder_present;
    const char                *filter_title;
    H5Z_can_apply_func_t      can_apply;
    H5Z_set_local_func_t      set_local;
    H5Z_func_t                filter;
    H5Z_set_config_func_t     set_config;
    H5Z_get_config_func_t     get_config;
} H5Z_class3_t;
\end{minted}

\verysubsection{Field descriptions}

\begin{itemize}
\item \texttt{filter\_title}---An optional short human-readable label for the
  filter (e.g., \texttt{"ZFP Compression"}, \texttt{"Blosc2"}). This field is
  not used for registry lookup or any programmatic dispatch---it is purely
  informational. When non-\texttt{NULL}, \texttt{H5Zregister()} packs the value
  directly into \texttt{cd\_values} as the reserved key \texttt{filter\_title}
  after \texttt{set\_config} has populated the filter's own \texttt{cd\_values}.
  The key is never passed to \texttt{set\_config} or \texttt{get\_config}. Because
  it is packed into \texttt{cd\_values} alongside the filter's own configuration,
  it persists in the file and is available to tools such as \texttt{h5dump} even
  when the plugin is not installed. The value must be printable ASCII with no
  leading or trailing whitespace and no equals signs. It may be \texttt{NULL}.
  \texttt{H5Pget\_filter2} returns \texttt{filter\_title} (extracted from
  \texttt{cd\_values}) when present, and falls back to the numeric filter ID
  formatted as a decimal string otherwise.

  \textbf{Length limit:} \texttt{filter\_title} must not exceed 255 bytes
  (excluding the NUL terminator). \texttt{H5Zregister()} rejects a v3
  registration where \texttt{strlen(filter\_title) > 255} with
  \texttt{H5E\_BADVALUE}.  This bound guarantees that the packed
  \texttt{filter\_title} key consumes at most
  $\lceil 255/4 \rceil + 1 = 65$ \texttt{cd\_values} slots, leaving the
  remaining capacity exclusively for the filter's own parameters.

  \textbf{Legacy plugin hazard:} Because \texttt{filter\_title} is appended
  to \texttt{cd\_values}, a file written by a v3-aware library and then
  opened by an older library (e.g., HDF5~1.14.x) with a v2 plugin will
  present the older plugin's execution callback with a \texttt{cd\_nelmts}
  value larger than the plugin expects.  Plugins that strictly validate
  \texttt{cd\_nelmts} will return an error, breaking read-compatibility for
  those datasets.

  \textbf{Mitigation:} Plugin authors are strongly encouraged \emph{not} to
  validate \texttt{cd\_nelmts} as an exact equality check; the HDF5
  developer documentation already advises treating unknown trailing
  \texttt{cd\_values} slots as reserved.  New built-in filter callbacks
  upgraded to v3 explicitly tolerate extra trailing slots.  Plugin authors
  who need strict validation should check only that \texttt{cd\_nelmts}
  is \emph{at least} the expected count.  This guidance is documented in
  \texttt{H5Zdevelop.h} and in the plugin migration guide produced
  alongside this RFC.

  A migration guide for plugin authors upgrading from v2 to v3 is produced
  alongside this RFC, covering at minimum: the new \texttt{filter\_title} field
  semantics, the version constant change, and the locations in the library where
  v2/v3 branching is required.
\end{itemize}

\subsection{Callback Specifications}
\label{sec:callback-specs}

\verysubsection{\texttt{H5Z\_set\_config\_func\_t}}

\begin{itemize}
\item \texttt{params}---The key-value parameter string (e.g.,
  \texttt{"mode=rate, rate=3.5"}). \texttt{NULL} if no parameters were provided.
\item \texttt{flags}---On input, contains the flags value from the API call. The
  callback can modify this value; the modified value is what
  \texttt{H5Pappend\_filter} writes into the DCPL. This gives the filter the
  final say on properties like \texttt{H5Z\_FLAG\_MANDATORY} vs.\
  \texttt{H5Z\_FLAG\_OPTIONAL}. Callers that need to detect a flag override can
  compare the value they passed with the flags returned by
  \texttt{H5Pget\_filter2} after the call.
\item \texttt{cd\_nelmts}---On input, points to a value of~0. On output, set to the
  number of elements written to \texttt{cd\_values} (or the number of elements
  required, if \texttt{cd\_values} is \texttt{NULL}).
\item \texttt{cd\_values}---On output, populated with the filter's client data. If
  \texttt{NULL}, the callback still sets \texttt{cd\_nelmts} to the required count,
  enabling a two-pass size query pattern. The library calls the callback once
  with \texttt{cd\_values=NULL} to determine the required size, allocates
  accordingly, and calls it again with a sufficiently sized array.

  \textbf{Two-pass identity contract:} The callback returns the same
  \texttt{cd\_nelmts} value on both invocations. If the second-pass \texttt{cd\_nelmts}
  differs from the first-pass value in either direction, the library
  returns an error. A smaller second-pass value would silently truncate the
  populated array; a larger value would constitute a buffer overflow. Both
  represent a buggy callback. Document this contract in
  \texttt{H5Zdevelop.h} as part of the \texttt{H5Z\_set\_config\_func\_t} specification.

  Additionally, the library validates on the second pass that
  \texttt{cd\_nelmts} does not exceed the allocated capacity; if it does, the
  library returns an error rather than writing out of bounds.

  The library also verifies that \texttt{cd\_nelmts} does not exceed
  \texttt{UINT16\_MAX}, the current file-format limit on the number of
  \texttt{cd\_values} elements per filter entry.  A callback returning a count
  larger than this is rejected with \texttt{H5E\_BADVALUE}.  The parameter
  type is \texttt{size\_t} rather than a narrower type so that the limit can be
  raised in a future format revision without changing the callback signature.
  Note: the per-string token count is separately bounded by
  \texttt{H5Z\_CONFIG\_MAX\_PARAMS} (64), a parser-level limit applied before
  \texttt{set\_config} is ever invoked (\SectionRef{sec:param-overview}).
\item \texttt{cd\_values\_size}---The capacity of the \texttt{cd\_values} array in
  bytes (0 when \texttt{cd\_values} is \texttt{NULL}).
\item Returns non-negative on success, negative on failure.
\end{itemize}

The callback should use \texttt{H5Epush} to report meaningful error messages for
invalid parameters.

\verysubsection{Unknown-key handling}
A \texttt{set\_config} callback should reject unrecognized keys by returning a
negative value and pushing \texttt{H5E\_BADVALUE} with a message identifying the
key.  This is the primary defense against misspelled parameter names (e.g.,
\texttt{"clevl=6"} instead of \texttt{"clevel=6"}), which would otherwise be
silently ignored.  Because the string API does not have the compile-time safety
of named constants, rejecting unknown keys is strongly recommended for
usability.

However, a filter can choose to silently ignore unrecognized keys to support
forward compatibility: a caller that targets a newer version of a plugin can
write a parameter string containing new keys and have it accepted by an older
version of the same plugin without error, falling back to defaults for the
unknown parameters.  Filter authors who choose this approach should document it
clearly so that callers are not misled into believing all keys were acted upon.

All built-in filter \texttt{set\_config} callbacks reject unknown keys.  Document
this guidance in \texttt{H5Zdevelop.h} as part of the
\texttt{H5Z\_set\_config\_func\_t} contract.

\verysubsection{Locale requirement for \texttt{set\_config}}
\texttt{set\_config} callbacks that parse numeric values from the parameter
string (e.g., converting a rate or tolerance to \texttt{double}) use
C-locale numeric conversion. Process-wide locale settings (e.g., a
comma-decimal locale) do not affect parameter interpretation. In practice,
this means using \texttt{strtod} under \texttt{uselocale(newlocale(LC\_NUMERIC\_MASK,
"C", 0))} or a locale-independent equivalent. The library-provided helper
\texttt{H5Zconfig\_get\_param} (\SectionRef{sec:config-get-param}) returns raw
string values; it is the callback's responsibility to convert them to numeric
types in a locale-safe manner. This requirement mirrors the output-side locale
requirement on \texttt{get\_config} (\SectionRef{sec:callback-specs}) and ensures
round-trip correctness. Document this requirement in
\texttt{H5Zdevelop.h} as part of the \texttt{H5Z\_set\_config\_func\_t} specification.

\verysubsection{\texttt{H5Z\_get\_config\_func\_t}}

\begin{itemize}
\item Provides a best-effort human-readable representation of the filter
  configuration encoded in \texttt{cd\_values}.  Due to the lossy nature of the
  \texttt{cd\_values} encoding (e.g., floating-point values quantised to integers,
  information collapsed during the \texttt{set\_config} conversion), the output
  cannot in general be guaranteed to round-trip back to an identical
  configuration.  See the note on string reconstruction in
  \SectionRef{sec:introspection}.
\item Follows the standard HDF5 buffer-size query pattern: \texttt{buf\_size} is
  an in/out parameter. On entry, \texttt{*buf\_size} is the capacity of \texttt{buf};
  on return, \texttt{*buf\_size} is set to the number of bytes required (excluding
  the NUL terminator). If \texttt{buf} is \texttt{NULL} or the capacity is~0, the
  callback still sets \texttt{*buf\_size} to the required size and returns
  success, enabling a two-pass size query.
\item Returns \texttt{herr\_t}: zero on success, negative on failure.
\item The output conforms to the grammar defined in
  \SectionRef{sec:param-overview} so that it is at least valid input to
  \texttt{set\_config}, even if the resulting configuration may differ from the
  original due to precision loss or encoding limitations.
\item \textbf{Locale requirement:} \texttt{get\_config} implementations that format
  numeric values (e.g., floating-point parameters) use the \texttt{"C"} locale
  for number formatting. Process-wide locale settings (e.g., a locale that
  uses a comma as the decimal separator) do not affect the output. In
  practice, this means using \texttt{snprintf} with \texttt{uselocale(newlocale(LC\_NUMERIC\_MASK,
  "C", 0))} or equivalent, or formatting numeric values through a
  locale-independent method. Outputting \texttt{"rate=3,5"} in a
  comma-decimal locale would violate the parameter string grammar and break
  round-tripping via \texttt{set\_config}. Document this requirement
  in \texttt{H5Zdevelop.h} as part of the \texttt{H5Z\_get\_config\_func\_t} specification.
\end{itemize}


\subsection{Backward Compatibility}

\texttt{H5Zregister()} inspects the \texttt{version} field to determine whether the
structure is \texttt{H5Z\_class1\_t} (legacy), \texttt{H5Z\_class2\_t}, or
\texttt{H5Z\_class3\_t}. Existing v1 and v2 plugins continue to function without
modification. The new callbacks and the \texttt{filter\_title} field are never
accessed for v1/v2 registrations.
\texttt{H5Z\_class3\_t} is a new type, not a modification of \texttt{H5Z\_class2\_t}.

A v3 plugin remains fully usable through the traditional \texttt{H5Pset\_filter} API.
Upgrading to v3 is optional and additive.

Files written with a \texttt{filter\_title} key in \texttt{cd\_values} are
readable by older HDF5 libraries through \emph{existing} v2 plugins, provided
those plugins tolerate a \texttt{cd\_nelmts} value larger than their baseline
expectation (see the legacy plugin hazard note in \SectionRef{sec:class3-struct}).
Plugin authors who perform exact \texttt{cd\_nelmts} validation must update their
callbacks to accept trailing slots before their plugin is used with a v3-aware
library.

For v3 plugins, if \texttt{filter\_title} is non-\texttt{NULL}, after
\texttt{set\_config} has populated the filter's \texttt{cd\_values}, the library
appends a \texttt{filter\_title=\textit{value}} slot directly into \texttt{cd\_values}.
This key is never seen by \texttt{set\_config} or \texttt{get\_config}. When
\texttt{H5Pget\_filter2} is called, it extracts \texttt{filter\_title} from
\texttt{cd\_values} if present, and falls back to the numeric filter ID formatted
as a decimal string otherwise.


%% ======================================================================
%%  Section 5 – Public API
%% ======================================================================
\section{Public API}
\label{sec:public-api}

\subsection{Single-Filter Configuration}
\label{sec:append-filter}

\verysubsection{Parameter type}

\begin{minted}{c}
typedef enum {
    H5Z_PARAMS_CDVALUES = 0,  /* raw cd_values array (same as H5Pset_filter) */
    H5Z_PARAMS_STRING   = 1,  /* human-readable key=value string              */
} H5Z_params_type_t;

typedef struct {
    H5Z_params_type_t type;
    union {
        struct {
            size_t          cd_nelmts;
            const unsigned *cd_values;
        } raw;
        const char *str;
    } u;
} H5Z_params_t;
\end{minted}

Convenience initialiser macros are provided for the two common cases:

\begin{minted}{c}
/* C99: compound literals with designated initialisers */
#define H5Z_PARAMS_RAW(n, vals) \
    ((H5Z_params_t){ H5Z_PARAMS_CDVALUES, { .raw = { (n), (vals) } } })
#define H5Z_PARAMS_STR(s) \
    ((H5Z_params_t){ H5Z_PARAMS_STRING,   { .str  = (s)            } })
\end{minted}

\verysubsection{Compound literal lifetime}
Both macros expand to C99 compound literals.  A compound literal has
automatic storage duration within its enclosing block---it is not a
temporary that expires at the semicolon.  Passing \texttt{\&H5Z\_PARAMS\_STR("level=6")}
to \texttt{H5Pappend\_filter} is therefore valid C99: the literal's
address is well-defined for the duration of the call.  However, storing
that address in a pointer that outlives the enclosing block is undefined
behaviour; callers who need to store an \texttt{H5Z\_params\_t} beyond a
single call should declare a named variable instead:

\begin{minted}{c}
/* Safe: compound literal address taken for the duration of one call */
H5Pappend_filter(plist, H5Z_FILTER_DEFLATE,
                     H5Z_FLAG_MANDATORY, &H5Z_PARAMS_STR("level = 6"));

/* Also safe: named variable with explicit lifetime */
H5Z_params_t p = H5Z_PARAMS_STR("level = 6");
H5Pappend_filter(plist, H5Z_FILTER_DEFLATE, H5Z_FLAG_MANDATORY, &p);
\end{minted}

\verysubsection{Function signature}

\begin{minted}{c}
herr_t H5Pappend_filter(
    hid_t              plist_id,
    H5Z_filter_t       id,
    unsigned           flags,
    const H5Z_params_t *params
);
\end{minted}

\verysubsection{Description}
Configures the filter identified by \texttt{id} and appends it to the
existing filter pipeline on \texttt{plist\_id}.  The pipeline is subject to
the existing maximum of \texttt{H5Z\_MAX\_NFILTERS} (32) entries.

\texttt{H5Pappend\_filter} accepts either a raw \texttt{cd\_values} array
(\texttt{H5Z\_PARAMS\_CDVALUES}) or a key=value parameter string
(\texttt{H5Z\_PARAMS\_STRING}) through the tagged union \texttt{H5Z\_params\_t}.
Passing \texttt{NULL} for \texttt{params} is equivalent to
\texttt{H5Z\_PARAMS\_CDVALUES} with \texttt{cd\_nelmts = 0} and is valid for
filters that take no parameters.

The name \texttt{H5Pappend\_filter} reflects its semantics: each call appends
a new entry to the filter pipeline.  Modifying the parameters of an existing
pipeline entry is handled by \texttt{H5Pmodify\_filter} (see the out-of-scope
note in \SectionRef{sec:intro}).

When \texttt{type} is \texttt{H5Z\_PARAMS\_CDVALUES}, the function behaves
identically to \texttt{H5Pset\_filter}: the \texttt{cd\_values} array is stored
in the DCPL without invoking any plugin callback.

When \texttt{type} is \texttt{H5Z\_PARAMS\_STRING}, the library locates or
loads the filter plugin, invokes the filter's \texttt{set\_config} callback to
translate the parameter string into \texttt{cd\_values}, and stores the result.

\textbf{Empty-string fast path:} If \texttt{u.str} is \texttt{NULL} or the empty
string \texttt{""}, \texttt{H5Pappend\_filter} treats the call as ``no
parameters''---equivalent to \texttt{H5Z\_PARAMS\_CDVALUES} with
\texttt{cd\_nelmts~=~0}---regardless of whether the filter implements
\texttt{set\_config}. The plugin is not loaded and no callback is invoked.
This makes parameterless filters (shuffle, fletcher32, nbit) usable through
the \texttt{H5Z\_PARAMS\_STRING} form without forcing them to provide a
\texttt{set\_config} callback that has nothing to do.

\verysubsection{Eager plugin loading}
\textbf{Note:} Unlike \texttt{H5Pset\_filter}, which stores raw \texttt{cd\_values}
without consulting the plugin, \texttt{H5Pappend\_filter} with
\texttt{H5Z\_PARAMS\_STRING} loads the filter plugin and invokes its
\texttt{set\_config} callback at call time. The plugin must therefore be
available when \texttt{H5Pappend\_filter} is called, not merely when
\texttt{H5Dcreate} runs. If the plugin cannot be found,
\texttt{H5Pappend\_filter} returns \texttt{H5E\_NOFILTER} immediately rather
than deferring the failure to \texttt{H5Dcreate}. This is a deliberate
design choice---fail-fast at configuration time gives clearer errors than
deferred failure during dataset creation---but it differs from the lazy
behavior of \texttt{H5Pset\_filter} and may surprise callers porting
existing code.

\verysubsection{Behavioral note on flags}
The caller controls flags directly.  When \texttt{type} is
\texttt{H5Z\_PARAMS\_STRING}, the flags value is passed through the filter's
\texttt{set\_config} callback, which may modify it (e.g., SZIP unconditionally
forces \texttt{H5Z\_FLAG\_OPTIONAL}).  Callers that need to detect a flag
override can compare the value they passed with the flags returned by
\texttt{H5Pget\_filter2} after the call.  See \SectionRef{sec:callback-specs}.

\verysubsection{Parameters}

\begin{itemize}
\item \texttt{plist\_id}---A dataset creation property list identifier.
\item \texttt{id}---The filter identifier (\texttt{H5Z\_filter\_t}).  For
  built-in filters, use the predefined constants
  (\texttt{H5Z\_FILTER\_DEFLATE}, etc.).  For externally registered filters,
  use the numeric ID assigned by the HDF Group filter registry.
\item \texttt{flags}---A bit mask specifying filter behavior
  (\texttt{H5Z\_FLAG\_MANDATORY}, \texttt{H5Z\_FLAG\_OPTIONAL}).
\item \texttt{params}---A pointer to an \texttt{H5Z\_params\_t} describing the
  filter configuration, or \texttt{NULL} for filters that take no parameters.
  \begin{itemize}
  \item \texttt{H5Z\_PARAMS\_CDVALUES}: \texttt{u.raw.cd\_nelmts} and
    \texttt{u.raw.cd\_values} are used directly, identical to
    \texttt{H5Pset\_filter}.
  \item \texttt{H5Z\_PARAMS\_STRING}: \texttt{u.str} is a comma-separated
    key=value parameter string.  \texttt{NULL} or an empty string means ``no
    parameters'' and is valid for parameterless filters (shuffle, fletcher32,
    nbit).
  \end{itemize}
\end{itemize}

\verysubsection{Returns}
Non-negative on success; negative on failure.

\verysubsection{Errors}

\begin{itemize}
\item \texttt{H5E\_NOFILTER}---(\texttt{H5Z\_PARAMS\_STRING} only) The filter
  plugin could not be found or loaded.
\item \texttt{H5E\_UNSUPPORTED}---(\texttt{H5Z\_PARAMS\_STRING} only) The
  filter was found but does not implement \texttt{set\_config}, and a
  non-empty parameter string was supplied.  An empty string or
  \texttt{NULL} is accepted regardless of whether \texttt{set\_config} is
  present (see the empty-string fast path in the Description).
\item \texttt{H5E\_BADVALUE}---The parameter string is malformed, the
  \texttt{params} type field is unrecognised, or \texttt{set\_config} rejected
  the parameter string.
\item \texttt{H5E\_CANTINIT}---The \texttt{set\_config} callback returned a
  \texttt{cd\_nelmts} value exceeding \texttt{UINT16\_MAX} or the library's
  internal filter-table limit.
\end{itemize}

\verysubsection{Usage examples}

\begin{minted}{c}
/* --- String form (TOML inline table content) --- */

/* Built-in deflate: integer value */
H5Pappend_filter(plist, H5Z_FILTER_DEFLATE, H5Z_FLAG_MANDATORY,
                     &H5Z_PARAMS_STR("level = 6"));

/* External plugin (ZFP): string + float values */
H5Pappend_filter(plist, H5Z_FILTER_ZFP, H5Z_FLAG_MANDATORY,
                     &H5Z_PARAMS_STR("mode = \"rate\", rate = 3.5"));

/* Boolean flag */
H5Pappend_filter(plist, H5Z_FILTER_BLOSC, H5Z_FLAG_MANDATORY,
                     &H5Z_PARAMS_STR("fast_mode = true, level = 6"));

/* Parameterless filter */
H5Pappend_filter(plist, H5Z_FILTER_SHUFFLE, H5Z_FLAG_MANDATORY, NULL);

/* --- Raw cd_values form (equivalent to H5Pset_filter) --- */

unsigned level = 6;
H5Pappend_filter(plist, H5Z_FILTER_DEFLATE, H5Z_FLAG_MANDATORY,
                     &H5Z_PARAMS_RAW(1, &level));

/* --- Building a mixed pipeline --- */
H5Pappend_filter(plist, H5Z_FILTER_SHUFFLE,    H5Z_FLAG_MANDATORY, NULL);
H5Pappend_filter(plist, H5Z_FILTER_DEFLATE,    H5Z_FLAG_MANDATORY,
                     &H5Z_PARAMS_STR("level = 9"));
H5Pappend_filter(plist, H5Z_FILTER_FLETCHER32, H5Z_FLAG_MANDATORY, NULL);
\end{minted}


\subsection{Pipeline Introspection}
\label{sec:introspection}

\begin{minted}{c}
herr_t H5Pget_filter_params_by_idx(
    hid_t     plist_id,
    unsigned  filter_idx,
    char     *params_buf,
    size_t    params_buf_size,
    size_t   *params_len
);
\end{minted}

\verysubsection{Description}
Returns the human-readable parameter string for the filter at position
\texttt{filter\_idx} in the pipeline. For filters that implement \texttt{get\_config},
the callback's output is used. For filters that do not implement
\texttt{get\_config}, the raw \texttt{cd\_values} are formatted as a fallback string
using colon separators (e.g., \texttt{"cd\_values=1:0:0:1074921472"}). The colon
separator avoids ambiguity with the comma-separated key=value grammar defined in
\SectionRef{sec:param-format}---the fallback is a single key=value pair whose
value contains colons, which is valid under the grammar.

\verysubsection{Note on string reconstruction}
The DCPL does not retain the original parameter string. After
\texttt{H5Pappend\_filter} resolves the string into \texttt{cd\_values}, only
the \texttt{cd\_values} are stored in the property list (identical to
\texttt{H5Pset\_filter}). This function relies on the filter's \texttt{get\_config}
callback to reconstruct a human-readable string from the stored \texttt{cd\_values}.

This reconstruction is inherently limited. The \texttt{cd\_values} encoding was
designed to carry opaque integer data, not to serve as a lossless serialisation
of a key-value string. For some filters (e.g., ZFP), floating-point
configuration values are quantised during conversion to unsigned integers; the
original value cannot be exactly recovered from the stored integer, and the
reconstructed string will reflect the rounded value rather than what the user
originally typed. Consumers should not rely on exact string equality between
input and output, and should treat the reconstructed string as informational
only.

\verysubsection{Note on \texttt{H5Pget\_filter2}}
The existing \texttt{H5Pget\_filter2} function returns a filter name from the class
struct's \texttt{name} field (the ``debug comment'' in v2 plugins). For v3 plugins,
\texttt{H5Pget\_filter2} extracts \texttt{filter\_title} from \texttt{cd\_values} if
the key is present, and falls back to the numeric filter ID formatted as a decimal
string otherwise. Because \texttt{filter\_title} is packed into \texttt{cd\_values}
at registration time, this human-readable label is available even when the plugin
is not loaded.
\texttt{H5Pget\_filter\_params\_by\_idx} complements \texttt{H5Pget\_filter2} by returning
the parameter string rather than the filter title.

\verysubsection{Parameters}

\begin{itemize}
\item \texttt{plist\_id}---A dataset creation property list identifier.
\item \texttt{filter\_idx}---Zero-based index of the filter in the pipeline.
\item \texttt{params\_buf}---Buffer to receive the parameter string, or \texttt{NULL}
  to query size only.
\item \texttt{params\_buf\_size}---Size of \texttt{params\_buf} in bytes.
\item \texttt{params\_len}---On output, set to the length of the parameter string
  (excluding NUL), or \texttt{NULL} if not needed. Set even when \texttt{params\_buf}
  is \texttt{NULL}, enabling a size query.
\end{itemize}

\verysubsection{Returns}
Non-negative on success; negative on failure.


\subsection{Developer Helpers: Typed Parameter Lookup}
\label{sec:config-get-param}

Because the parameter string format is now TOML-typed, the library provides
typed accessor functions in \texttt{H5Zdevelop.h} rather than a single
string-returning helper.  All built-in filter \texttt{set\_config} implementations
use these functions.

\begin{minted}{c}
/* Key presence check (all types) */
htri_t H5Zconfig_has_key(const char *params, const char *key);

/* Typed accessors -- return htri_t: >0 found, 0 absent, <0 error */
htri_t H5Zconfig_get_int    (const char *params, const char *key,
                              int64_t  *out);
htri_t H5Zconfig_get_double (const char *params, const char *key,
                              double   *out);
htri_t H5Zconfig_get_bool   (const char *params, const char *key,
                              hbool_t  *out);
htri_t H5Zconfig_get_str    (const char *params, const char *key,
                              char *buf, size_t *buf_size);
\end{minted}

\verysubsection{Description}
Each function searches the parameter string \texttt{params} for \texttt{key}
and writes the decoded value into \texttt{*out} (or \texttt{buf} for the string
accessor).  The \texttt{htri\_t} return type separates three cases cleanly:
found (\texttt{>~0}), absent (\texttt{0}), and error (\texttt{<~0}) ---
without pushing to the error stack for a merely absent optional parameter.

\texttt{H5Zconfig\_get\_str} follows the standard HDF5 size-query pattern:
\texttt{buf\_size} is in/out.  On entry \texttt{*buf\_size} is the buffer
capacity; on return it is set to the required byte count (excluding NUL).
Passing \texttt{NULL} for \texttt{buf} performs a size query only.  If the
buffer is too small, \texttt{H5E\_OVERFLOW} is pushed and the function returns
negative; the caller reallocates and retries.

Type mismatches (e.g., calling \texttt{H5Zconfig\_get\_int} on a key whose
value is a string) return negative with \texttt{H5E\_BADVALUE}.

These functions are placed in \texttt{H5Zdevelop.h} to limit the public API
surface.


\subsection{Normative \texttt{cd\_values} Packing Convention}
\label{sec:cd-packing-convention}

For \texttt{set\_config} and \texttt{get\_config} callbacks to round-trip
correctly across plugins and tools, typed parameters wider than a single
\texttt{unsigned int} must be packed into \texttt{cd\_values} using a
consistent layout.  The following convention is normative; built-in filter
callbacks that handle floating-point or string parameters use it, and
third-party filters are strongly encouraged to do the same.

\begin{center}
\renewcommand{\arraystretch}{1.2}
\begin{tabularx}{\textwidth}{llX}
\toprule
\textbf{C type} & \textbf{Slots} & \textbf{Encoding} \\
\midrule
\texttt{unsigned int} & 1 & direct store \\
\texttt{int32\_t}     & 1 & two's-complement bitcast to \texttt{uint32} \\
\texttt{float}        & 1 & IEEE\,754 \texttt{memcpy} into one slot \\
\texttt{double}       & 2 & IEEE\,754 \texttt{memcpy}; slot[0]\,=\,low\,32\,bits (LE), slot[1]\,=\,high\,32\,bits \\
\texttt{int64\_t}     & 2 & two's-complement \texttt{memcpy}; slot[0]\,=\,low\,32\,bits (LE), slot[1]\,=\,high\,32\,bits \\
\texttt{char[]}       & $1{+}\lceil\text{len}/4\rceil$ &
                          slot[0]\,=\,byte length; remaining slots hold 4-byte null-padded chunks \\
\bottomrule
\end{tabularx}
\end{center}

For \texttt{double} and \texttt{int64\_t} the two-slot layout is always
little-endian regardless of the host platform; on big-endian hosts the
plugin must perform the necessary byte swap, matching the existing HDF5
convention that all \texttt{cd\_values} are stored in little-endian order.
The convention assumes IEEE\,754 \texttt{float} (32~bits) and \texttt{double}
(64~bits); types not enumerated above (e.g., \texttt{long double}) are out
of scope for this encoding and must be reduced to one of the listed types
by the plugin.

The library does not export typed pack/unpack helpers as part of the public
API; plugin authors implement the convention directly using \texttt{memcpy}
and the byte layout above.  The same convention is documented in
\texttt{H5Zdevelop.h} alongside the \texttt{set\_config} specification.


\subsection{Canonical \texttt{set\_config} Implementation Template}
\label{sec:set-config-template}

The two-pass \texttt{set\_config} protocol (\SectionRef{sec:callback-specs}) is
subtle. Plugin authors ensure identical \texttt{cd\_nelmts} on both passes.
The following template shows the recommended implementation pattern. All
built-in filter \texttt{set\_config} callbacks follow this structure and are
documented in \texttt{H5Zdevelop.h} as normative examples.

\begin{minted}{c}
static herr_t
my_filter_set_config(const char *params, unsigned *flags,
                     size_t *cd_nelmts, unsigned cd_values[],
                     size_t cd_values_size)
{
    /* ---- 1. Parse parameters (identical on both passes) ---- */
    int64_t level = 6;   /* default */

    if (params) {
        htri_t found = H5Zconfig_get_int(params, "level", &level);
        if (found < 0) return FAIL;
        if (found > 0 && (level < 0 || level > 9)) {
            H5Epush(..., H5E_BADVALUE, "level must be 0-9");
            return FAIL;
        }
    }

    /* ---- 2. Report element count (SAME on both passes) ---- */
    *cd_nelmts = 1;

    /* ---- 3. Populate cd_values only on the second pass ---- */
    if (cd_values != NULL) {
        if (cd_values_size < *cd_nelmts * sizeof(*cd_values))
            return FAIL;               /* should not happen */
        cd_values[0] = (unsigned)level;
    }

    /* ---- 4. Optionally override flags ---- */
    /* (void)flags;  -- accept caller's flags unchanged */

    return SUCCEED;
}
\end{minted}

\verysubsection{Key invariants}
\begin{itemize}
\item \texttt{cd\_nelmts} is set to the same value regardless of whether
  \texttt{cd\_values} is \texttt{NULL} or non-\texttt{NULL}.
\item Parameter parsing is performed on both passes (not cached) to
  ensure the count is consistent.
\item If \texttt{cd\_values} is non-\texttt{NULL}, verify
  \texttt{cd\_values\_size >= *cd\_nelmts * sizeof(*cd\_values)} before writing.
\item Use the typed accessors (\texttt{H5Zconfig\_get\_int},
  \texttt{H5Zconfig\_get\_double}, etc.)\ rather than manual string conversion;
  locale independence and TOML type checking are handled by the accessors.
\end{itemize}

\textbf{Note on double parsing:} Because \texttt{cd\_nelmts} may vary depending on
the parameter values (not just their presence), the two-pass protocol requires
the callback to parse the parameter string twice.  Filters that wish to avoid
this overhead may report a fixed worst-case \texttt{cd\_nelmts} on both passes,
allocate that many slots on the stack, and fill unused trailing slots with a
sentinel value.  This is conforming provided the worst-case count is a genuine
constant---independent of the parameter string---so that both passes always
return the same value and the two-pass identity contract is satisfied.
The library does not mandate which approach is used.


\subsection{Pipeline-in-a-String API}
\label{sec:pipeline-string}

A convenience function that accepts a semicolon-delimited pipeline string
(e.g., \texttt{"shuffle; deflate, level = 9; fletcher32"}) is deferred to a
future release.  The semicolon character is reserved outside of quoted strings
(\SectionRef{sec:param-parser}) for this purpose.
\texttt{H5Pappend\_filter} (\SectionRef{sec:append-filter}) is the primary API.


%% ======================================================================
%%  Section – Current Limitations and Future On-Disk Format Work
%% ======================================================================
\section{Current Limitations and Future On-Disk Format Work}
\label{sec:pline-v3}

This RFC introduces \textbf{no on-disk format change}.  Files written via
\texttt{H5Pappend\_filter} are byte-identical to files written via
\texttt{H5Pset\_filter}: \texttt{cd\_values} are stored in the existing
\texttt{H5O\_PLINE\_VERSION\_2} object-header message.  The parameter string
is consumed at API call time and is not stored on disk.

\subsection{Limitations of the \texttt{cd\_values}-only Approach}

Without typed on-disk metadata, two limitations apply:

\begin{enumerate}
\item \textbf{Plugin-dependent introspection.}  Tools such as \texttt{h5dump}
  must load the filter plugin and call \texttt{get\_config} to produce a
  human-readable representation of \texttt{cd\_values}.  When the plugin is
  unavailable, the stored array appears as opaque integers.

\item \textbf{Best-effort round-trip.}  \texttt{get\_config} reconstructs a
  parameter string from \texttt{cd\_values}, but this reconstruction may be
  lossy for filters that store only a quantised approximation of a
  floating-point value.  Consumers must treat \texttt{get\_config} output as
  informational and must not rely on exact string equality with the original
  \texttt{H5Pappend\_filter} input.  Filters that use the
  \texttt{H5Z\_cd\_pack\_*} helpers (\SectionRef{sec:cd-packing-convention})
  can achieve lossless round-trips.
\end{enumerate}

These limitations are shared with the existing \texttt{H5Pset\_filter} API.

\subsection{Future Work: Typed On-Disk Storage}

A future RFC targeting HDF5~3.0 will propose a typed on-disk storage
extension---provisionally \texttt{H5O\_PLINE\_VERSION\_3}---that appends
per-slot type tags alongside \texttt{cd\_values}.  The
\texttt{H5Z\_cd\_pack\_*} packing convention defined in this RFC is the
normative foundation: filters that adopt it today will be forward-compatible
with that extension without changes to their callbacks.

The format extension is deferred because the primary value of this RFC---human-readable
filter configuration at the API level---is fully realised without a format
change, and requiring a new libver upper bound would slow adoption.


\end{document}
